\scanpage{406}\mspage{145}
9.11.90

Pour les questions précédentes, relatives aux flèches
dans le diagramme
\diagram{\begin{tikzcd}[ampersand replacement=\&,row sep=small,column sep=small]
\& X\arrow[d,"{\unsure{i|\tau}}"] \& \\
\& \underline{\mathrm{Ch}}_{\infty}^{\deleted{\unsure{t}}}(X)\arrow[d,"p"]\arrow[ddl,"\sigma"']\arrow[ddr,"\beta",-{Circle[length=1.2pt]}]\&\\
\&X\times X\arrow[dl,"{\unsure{p_1|r_1}}"']\arrow[dr,no head,"p_2"]\&\\
X\&\&\unsure{\bar X|X}
\end{tikzcd}}
% slanted
\margin{\deleted{\unsure{figure}}}
, il faut préciser absolument
la description de $\underline{\mathrm{Ch}}_{\infty}(X)$ \uline{dans
$\mathrm{Cat}$}, donc il faut préciser
le système projectif $\underline t\deleted{\unsure{_{\infty}}}=(t_n)_{n\in\mathbb N}$.
Nous allons prendre en effet
\[
t_n=(\underbrace{\deleted{\unsure{\relbar}}\Delta_+\Delta_-}_{\tau_0})^{2n},
\]
avec les morphismes de
transition décrits par ailleurs
\diagram{\begin{tikzcd}[ampersand replacement=\&,row sep=small]
t_{n+p}\arrow[r,{Bar[width=2pt]}-{Bar[width=2pt]}]\arrow[d,phantom,sloped,"=" description]\&t_n\arrow[d,phantom,sloped,"=" description]\\
\tau_0^{2n+2p}\&\tau_0^{\unsure{n|2n}}
\end{tikzcd}}
étant donné par $\delta_p*\mathrm{id}_{\tau_0^{2n}}*\delta_p$, où $\delta_p\colon\tau_0^p\unsure{\relbar|\to}1$ est
la dégénérescence de $\tau_0^p$. Ainsi, on peut
considérer $\mathrm{Ch}_{\infty}(X)$ comme étant la catégorie
des \deleted{\unsure{entre}} chemins \inserted{bi-}infinis dans $X$
\diagram{\begin{tikzcd}[ampersand replacement=\&,column sep=1.8pt,font=\scriptsize,execute at end picture={
\draw (bl.north) .. controls +(0,.44) and +(-.32,0) .. ($(az.north)+(0,.72)$) .. controls +(.32,-.08) and +(0,.44) .. (br.north);
\node[above] at ($(az.north)+(0,.72)$) {$\underline{\mathrm{Ch}}_2$};
\draw (al.north) .. controls +(0,.27) and +(-.28,-.13) .. ($(az.north)+(0,.40)$) .. controls +(.28,-.13) and +(0,.27) .. (ar.north);
\node[above] at ($(az.north)+(0,.40)$) {$\underline{\mathrm{Ch}}_1$};
\node[above,inner sep=0pt] at (az.north) {$\underset{\downarrow}{\underline{\mathrm{Ch}}_0}$};}]
\cdots\arrow[r]\&a_{-5}\&|[alias=bl]|a_{-4}\arrow[l]\arrow[r]\&a_{-3}\&|[alias=al]|a_{-2}\arrow[l,"u_{-2}"']\arrow[r,"u_{-1}"]\&a_{-1}\&|[alias=az]|a_0\arrow[l,"u_0"']\arrow[r,"u_1"]\&a_1\&|[alias=ar]|a_2\arrow[l,"u_2"']\arrow[r,"u_3"]\&a_3\&|[alias=br]|a_4\arrow[l]\arrow[r]\&a_5\&\cdots\arrow[l]
\end{tikzcd}}
avec les $u_i$ des identités pour $|i|$ assez grand.
La sous-catégorie \deleted{\unsure{$\mathrm{Ch}_{t}$}} (pleine)
\deleted{Nous allons regarder cette catégorie}
\[
\underline{\mathrm{Ch}}^{W}_{t_n}(X)=\underline{\mathrm{Ch}}^{W}_{(\Delta_+\Delta_-)^{2n}}
\]
s'identifie à celle pour laquelle on a
\[
u_i\ \text{une identité pour }i\notin[-2\unsure{i|i'}+1,2i]
\]
i.e. les seuls $u_i$ ``significatifs'' sont les
$4i$ flèches $u_{-2i+1},u_{-2i+2}\ \ldots\ u_{2i}$.

Mais il sera \inserted{plus} commode de regarder $\underline{\mathrm{Ch}}_{\infty}(X)$
comme \deleted{\unsure{la}} limite inductive d'une
% archive sheet number: 405
